\chapter*{INTRODUCCIÓN}\addcontentsline{toc}{chapter}{INTRODUCCIÓN}\label{introduccion}

La congestión vehicular en Lima Metropolitana es uno de los problemas urbanos más críticos que afecta
diariamente a millones de personas. Un informe publicado por El Comercio en 2024 ubica a Lima entre las
ciudades con mayor nivel de congestión vehicular en América Latina, con una velocidad promedio cercana a
14.5 km/h. Este problema no solo representa una pérdida significativa de tiempo para los conductores, sino
que también genera un impacto económico relevante asociado a horas-hombre improductivas, mayor consumo de
combustible y menor eficiencia logística. Por otra parte, el transporte contribuye de manera importante a
las emisiones contaminantes metropolitanas (Banco Mundial, 2024), por lo que la mejora del control del
tránsito también tiene implicancias ambientales.

Una de las principales causas de esta congestión es la ineficiencia del sistema semafórico actual, que en
gran medida opera con ciclos fijos. En Lima Metropolitana existen cerca de 1\,200 intersecciones
semaforizadas bajo administración municipal, de las cuales una amplia mayoría opera de forma descoordinada
y sin gestión centralizada (Diario Correo, 2020). La mayor parte de los semáforos no tiene capacidad de
adaptarse a las variaciones del flujo vehicular en tiempo real; solo un conjunto reducido de intersecciones
cuenta con semáforos denominados «inteligentes», mientras que el resto opera de manera aislada y sin
conexión a un centro de control. Este modelo rígido aumenta los tiempos de espera, da lugar a una
utilización ineficiente de recursos como el combustible y el tiempo de los conductores, e incrementa los
niveles de contaminación (Sacyr, 2021).

Frente a este escenario, la modernización de los sistemas semafóricos mediante tecnologías inteligentes es
una necesidad para optimizar la movilidad urbana y reducir sus impactos negativos. Sin embargo, la
sustitución completa de la infraestructura instalada resulta costosa y de implementación lenta. La presente
tesis aborda el problema como un desafío de integración mecatrónica: el sistema debe observar el estado de
la intersección, transformar mediciones en variables de tráfico, calcular ajustes dentro de límites
operativos, conservar capacidad de decisión local ante pérdida de conectividad y comunicar telemetría sin
comprometer la seguridad funcional del controlador vial. Para ello se diseña un módulo IoT complementario
con visión por computadora, procesamiento en el borde, control difuso y una plataforma en la nube destinada
a supervisión, almacenamiento histórico y distribución controlada de parámetros.

El diseño se sustenta en las lecciones aprendidas de países con sistemas semafóricos avanzados y de
experiencias exitosas en Latinoamérica, adaptando las mejores prácticas a las necesidades locales. El
documento se organiza en siete capítulos: antecedentes y formulación del problema; estado del arte; diseño
conceptual; diseño electrónico; diseño mecánico; diseño de la etapa de control; e integración, validación
en SUMO y análisis de costos. Finalmente se presentan las conclusiones y recomendaciones derivadas del
trabajo.
